\subsection{Ambient sector retractions and normalized embeddings}

\begin{definition}[Ambient sector retractions and normalized embeddings]
    \label{def:mpdo_bnt_ambient_sector_maps}
    \lean{MPOTensor.BNTAlgebraTensorClause.oneSiteAmbientSectorRetraction}
    \lean{MPOTensor.BNTAlgebraTensorClause.oneSiteAmbientSectorRetraction_apply}
    \lean{MPOTensor.BNTAlgebraTensorClause.oneSiteNormalizedAmbientSectorRestoration}
    \lean{MPOTensor.BNTAlgebraTensorClause.oneSiteNormalizedAmbientSectorRestoration_apply}
    \lean{MPOTensor.BNTAlgebraTensorClause.TwoSiteMultiplicitySpectrum.%
        twoSiteAmbientSectorRetraction}
    \lean{MPOTensor.BNTAlgebraTensorClause.TwoSiteMultiplicitySpectrum.%
        twoSiteAmbientSectorRetraction_apply}
    \lean{MPOTensor.BNTAlgebraTensorClause.TwoSiteMultiplicitySpectrum.%
        twoSiteNormalizedAmbientSectorRestoration}
    \lean{MPOTensor.BNTAlgebraTensorClause.TwoSiteMultiplicitySpectrum.%
        twoSiteNormalizedAmbientSectorRestoration_apply}
    \lean{MPOTensor.BNTAlgebraTensorClause.TwoSiteMultiplicitySpectrum.%
        twoSiteNormalizedAmbientSectorRestoration_apply_explicit}
    \lean{MPOTensor.BNTAlgebraTensorClause.TwoSiteMultiplicitySpectrum.%
        relabeledTwoSiteMultiplicityTrace_eq_oneSite}
    \lean{MPOTensor.verticalSectorAmbientRetraction}
    \lean{MPOTensor.verticalSectorAmbientRetraction_apply}
    \lean{MPOTensor.verticalSectorAmbientRetraction_apply_sector}
    \lean{MPOTensor.normalizedVerticalSectorAmbientRestoration}
    \lean{MPOTensor.normalizedVerticalSectorAmbientRestoration_apply}
    \leanok
    \uses{def:mpdo_bnt_two_site_multiplicity_spectrum,
        def:mpdo_normalized_vertical_sector_maps,
        def:mpdo_retained_vertical_sector_coordinates,
        def:mpdo_two_site_blocking,
        thm:mpdo_bnt_multiplicity_trace_normalization}
    Let
    \begin{align}
        \mathcal A_1
        &=\bigoplus_\gamma \MN{d_\gamma},
        &
        \mathcal A_2^\sigma
        &=\bigoplus_\gamma \MN{\widetilde d_{\sigma(\gamma)}}
        \notag
    \end{align}
    be the one-site sector algebra and the relabelled two-site sector algebra.
    Their retained physical spaces are
    \begin{align}
        \mathcal K_1
        &=\bigoplus_\gamma
          \left(\C^{r_\gamma}\otimes\C^{d_\gamma}\right),
        &
        \mathcal K_2^\sigma
        &=\bigoplus_\gamma
          \left(\C^{\widetilde r_{\sigma(\gamma)}}
            \otimes\C^{\widetilde d_{\sigma(\gamma)}}\right).
        \notag
    \end{align}
    Let $W_1\colon\C^d\to\mathcal K_1$ and
    $W_2\colon\C^{d^2}\to\mathcal K_2^\sigma$ be the corresponding
    coisometries, so that $W_i W_i^\dagger=I$.  Write $\Pi_{1,\gamma}$ for
    the projection onto the $\gamma$-summand of $\mathcal K_1$ and
    $\Pi_{2,\gamma}^\sigma$ for the projection onto the
    $\sigma(\gamma)$-summand of $\mathcal K_2^\sigma$.  Let
    \begin{align}
        \kappa_\ast\colon\MN{d^2}
        &\longrightarrow \MN{d}\otimes\MN{d}
        \notag
    \end{align}
    be the canonical relabelling of a pair of physical indices.  For the
    positive diagonal multiplicity matrices $\mu_\gamma$ and
    $\nu_{\sigma(\gamma)}$, set
    \begin{align}
        m_\gamma
        &=\tr(\mu_\gamma)
         =\tr(\nu_{\sigma(\gamma)}).
        \label{eq:mpdo_bnt_ambient_sector_trace_weights}
    \end{align}
    Define
    \begin{align}
        R_1\colon\MN{d}&\longrightarrow\mathcal A_1,
        &
        [R_1(X)]_\gamma
        &=\tr_{\mathrm{mult}}
          \left(\Pi_{1,\gamma}W_1XW_1^\dagger\Pi_{1,\gamma}\right),
        \label{eq:mpdo_bnt_ambient_sector_R1}
        \\
        R_2\colon\MN{d}\otimes\MN{d}&\longrightarrow\mathcal A_2^\sigma,
        &
        [R_2(Y)]_\gamma
        &=\tr_{\mathrm{mult}}
          \left(\Pi_{2,\gamma}^\sigma W_2\kappa_\ast^{-1}(Y)W_2^\dagger
            \Pi_{2,\gamma}^\sigma\right),
        \label{eq:mpdo_bnt_ambient_sector_R2}
        \\
        \widetilde R_1\colon\mathcal A_1&\longrightarrow\MN{d},
        &
        \widetilde R_1((X_\gamma)_\gamma)
        &=W_1^\dagger
          \left(\bigoplus_\gamma
            \frac{\mu_\gamma}{m_\gamma}\otimes X_\gamma\right)W_1,
        \label{eq:mpdo_bnt_ambient_sector_R1_tilde}
        \\
        \widetilde R_2\colon\mathcal A_2^\sigma
          &\longrightarrow\MN{d}\otimes\MN{d},
        &
        \widetilde R_2((Y_\gamma)_\gamma)
        &=\kappa_\ast\!\left[
          W_2^\dagger\left(\bigoplus_\gamma
            \frac{\nu_{\sigma(\gamma)}}{m_\gamma}\otimes Y_\gamma\right)W_2
          \right].
        \label{eq:mpdo_bnt_ambient_sector_R2_tilde}
    \end{align}
    In Appendix~C.4 of \cite{Cirac2016MPDO_arXiv}, lines~2067--2068 define
    $R_1$ by~(\ref{eq:mpdo_bnt_ambient_sector_R1}), and lines~2078--2079
    define $R_2$ by~(\ref{eq:mpdo_bnt_ambient_sector_R2}).
    Lines~2081--2083 give the normalized map $\widetilde R_1$ in
    (\ref{eq:mpdo_bnt_ambient_sector_R1_tilde}).  Lines~2058--2071 give
    $\widetilde R_2$, including its common denominator, in
    (\ref{eq:mpdo_bnt_ambient_sector_R2_tilde}).
    These four definitions use only the two vertical decompositions, their
    sector relabelling, and
    (\ref{eq:mpdo_bnt_ambient_sector_trace_weights}).
\end{definition}

\begin{theorem}[Complete positivity of the ambient sector maps]
    \label{thm:mpdo_bnt_ambient_sector_complete_positivity}
    \lean{MPOTensor.BNTAlgebraTensorClause.oneSiteAmbientSectorRetractionExtension_isKrausCP}
    \lean{MPOTensor.BNTAlgebraTensorClause.oneSiteAmbientSectorRetractionExtension}
    \lean{MPOTensor.BNTAlgebraTensorClause.%
        oneSiteNormalizedAmbientSectorRestorationExtension_isKrausCP}
    \lean{MPOTensor.BNTAlgebraTensorClause.oneSiteNormalizedAmbientSectorRestorationExtension}
    \lean{MPOTensor.BNTAlgebraTensorClause.TwoSiteMultiplicitySpectrum.%
        twoSiteAmbientSectorRetractionExtension_isKrausCP}
    \lean{MPOTensor.BNTAlgebraTensorClause.TwoSiteMultiplicitySpectrum.%
        twoSiteAmbientSectorRetractionExtension}
    \lean{MPOTensor.BNTAlgebraTensorClause.TwoSiteMultiplicitySpectrum.%
        twoSiteNormalizedAmbientSectorRestorationExtension_isKrausCP}
    \lean{MPOTensor.BNTAlgebraTensorClause.TwoSiteMultiplicitySpectrum.%
        twoSiteNormalizedAmbientSectorRestorationExtension}
    \lean{MPOTensor.verticalSectorAmbientRetractionExtension}
    \lean{MPOTensor.verticalSectorAmbientRetractionExtension_factorization}
    \lean{MPOTensor.verticalSectorAmbientRetractionExtension_isKrausCP}
    \lean{MPOTensor.normalizedVerticalSectorAmbientRestorationExtension}
    \lean{MPOTensor.normalizedVerticalSectorAmbientRestorationExtension_factorization}
    \lean{MPOTensor.normalizedVerticalSectorAmbientRestorationExtension_isKrausCP}
    \leanok
    \uses{def:kraus_map_between_direct_sums,
        def:mpdo_bnt_ambient_sector_maps}
    Represent a sector family by its block-diagonal matrix.  For a map whose
    domain is a sector algebra, first extract the diagonal sector blocks from
    an arbitrary matrix.  With these canonical full-matrix representatives,
    each of $R_1$, $R_2$, $\widetilde R_1$, and $\widetilde R_2$ is completely
    positive.
\end{theorem}

\begin{proof}\leanok
    \uses{thm:mpdo_vertical_sector_partial_trace_kraus,
        lem:mpdo_single_kraus_map_cp,
        thm:mpdo_equiv_reindex_cptp,
        lem:is_kraus_cp_comp}
    Each raw retraction is the composition of compression by $W_i$, a unitary
    relabelling of coordinates, sectorwise partial trace, and block-diagonal
    inclusion.  Each normalized embedding is the composition of diagonal-block
    extraction, preparation of the positive matrix
    $\mu_\gamma/m_\gamma$ or
    $\nu_{\sigma(\gamma)}/m_\gamma$, coordinate relabelling, and expansion by
    $W_i^\dagger$.  Every constituent has a Kraus representation, and
    compositions of such maps are completely positive.
\end{proof}

% Local fix (zero-sector complement): the coisometries need not have full
% ambient support, so the raw maps are trace-nonincreasing rather than
% trace-preserving. Documented in
% docs/paper-gaps/cpgsv17_vertical_isometry_zero_sector.tex.
\begin{theorem}[Trace of the ambient sector maps]
    \label{thm:mpdo_bnt_ambient_sector_trace}
    \lean{MPOTensor.BNTAlgebraTensorClause.verticalSectorTrace_oneSiteAmbientSectorRetraction}
    \lean{MPOTensor.BNTAlgebraTensorClause.TwoSiteMultiplicitySpectrum.%
        verticalSectorTrace_twoSiteAmbientSectorRetraction}
    \lean{MPOTensor.BNTAlgebraTensorClause.verticalSectorTrace_oneSiteAmbientSectorRetraction_le}
    \lean{MPOTensor.BNTAlgebraTensorClause.TwoSiteMultiplicitySpectrum.%
        verticalSectorTrace_twoSiteAmbientSectorRetraction_le}
    \lean{MPOTensor.BNTAlgebraTensorClause.trace_oneSiteNormalizedAmbientSectorRestoration}
    \lean{MPOTensor.BNTAlgebraTensorClause.TwoSiteMultiplicitySpectrum.%
        trace_twoSiteNormalizedAmbientSectorRestoration}
    \lean{MPOTensor.BNTAlgebraTensorClause.%
        oneSiteNormalizedAmbientSectorRestorationExtension_isKrausCPTP}
    \lean{MPOTensor.BNTAlgebraTensorClause.TwoSiteMultiplicitySpectrum.%
        twoSiteNormalizedAmbientSectorRestorationExtension_isKrausCPTP}
    \lean{MPOTensor.verticalSectorTrace_verticalSectorAmbientRetraction_le}
    \lean{MPOTensor.trace_normalizedVerticalSectorAmbientRestoration}
    \lean{MPOTensor.normalizedVerticalSectorAmbientRestorationExtension_isKrausCPTP}
    \leanok
    \uses{def:mpdo_bnt_ambient_sector_maps,
        def:mpdo_positive_vertical_sector_family,
        thm:mpdo_bnt_ambient_sector_complete_positivity}
    For arbitrary ambient matrices $X\in\MN{d}$ and
    $Y\in\MN{d}\otimes\MN{d}$,
    \begin{align}
        \Trsec(R_1(X))
        &=\tr(W_1XW_1^\dagger),
        \label{eq:mpdo_bnt_ambient_sector_R1_trace}
        \\
        \Trsec(R_2(Y))
        &=\tr\!\left(W_2\kappa_\ast^{-1}(Y)W_2^\dagger\right).
        \label{eq:mpdo_bnt_ambient_sector_R2_trace}
    \end{align}
    If $X$ and $Y$ are positive semidefinite, then
    \begin{align}
        \Trsec(R_1(X))
        &\leq\tr(X),
        &
        \Trsec(R_2(Y))
        &\leq\tr(Y).
        \label{eq:mpdo_bnt_ambient_sector_trace_loss}
    \end{align}
    In the other direction, for $Z\in\mathcal A_1$ and
    $Z'\in\mathcal A_2^\sigma$,
    \begin{align}
        \tr(\widetilde R_1(Z))
        &=\Trsec(Z),
        &
        \tr(\widetilde R_2(Z'))
        &=\Trsec(Z').
        \label{eq:mpdo_bnt_ambient_sector_embedding_trace}
    \end{align}
    Hence the canonical full-matrix representatives of $\widetilde R_1$ and
    $\widetilde R_2$ are trace-preserving completely positive maps.

    The inequalities in
    (\ref{eq:mpdo_bnt_ambient_sector_trace_loss}) may be strict.  Indeed,
    $W_i^\dagger W_i$ is the projection onto the retained nonzero sectors and
    need not be the identity on the ambient space.  A positive matrix supported
    on the discarded orthogonal complement is annihilated by the corresponding
    compression.  Thus $R_1$ and $R_2$ are not asserted to preserve the trace
    on all ambient positive matrices, and no term acting on the discarded
    complement is included.
\end{theorem}

\begin{proof}\leanok
    \uses{lem:mpdo_normalized_vertical_sector_positive_trace,
        lem:mpdo_coisometric_compression_trace,
        thm:mpdo_equiv_reindex_cptp}
    Sectorwise partial trace and block-diagonal inclusion preserve the total
    trace, which gives
    (\ref{eq:mpdo_bnt_ambient_sector_R1_trace}) and
    (\ref{eq:mpdo_bnt_ambient_sector_R2_trace}).  For
    $P_i=W_i^\dagger W_i$, cyclicity of trace gives
    \begin{align}
        \tr(W_i ZW_i^\dagger)
        &=\tr(P_i Z)\leq\tr(Z)
        \notag
    \end{align}
    when $Z\succeq0$.  The relabelling $\kappa_\ast$ preserves positivity and
    trace, proving (\ref{eq:mpdo_bnt_ambient_sector_trace_loss}).  Finally, each
    normalized multiplicity matrix has trace one, and
    $W_i W_i^\dagger=I$ gives
    $\tr(W_i^\dagger BW_i)=\tr(B)$.  This proves
    (\ref{eq:mpdo_bnt_ambient_sector_embedding_trace}); complete positivity is
    supplied by the preceding theorem.
\end{proof}

\begin{theorem}[Retraction identities for the ambient sector maps]
    \label{thm:mpdo_bnt_ambient_sector_retractions}
    \lean{MPOTensor.BNTAlgebraTensorClause.%
        oneSiteAmbientSectorRetraction_comp_normalizedRestoration}
    \lean{MPOTensor.BNTAlgebraTensorClause.TwoSiteMultiplicitySpectrum.%
        twoSiteAmbientSectorRetraction_comp_normalizedRestoration}
    \lean{MPOTensor.verticalSectorAmbientRetraction_comp_normalizedRestoration}
    \leanok
    \uses{def:mpdo_bnt_ambient_sector_maps}
    The raw retractions are left inverses of their normalized embeddings:
    \begin{align}
        R_1\widetilde R_1
        &=\id_{\mathcal A_1},
        &
        R_2\widetilde R_2
        &=\id_{\mathcal A_2^\sigma}.
        \label{eq:mpdo_bnt_ambient_sector_retractions}
    \end{align}
    No identity is asserted for either reverse composite on the full ambient
    matrix algebra.
\end{theorem}

\begin{proof}\leanok
    \uses{lem:mpdo_retained_vertical_sector_retraction,
        thm:mpdo_equiv_reindex_cptp}
    Compressing an expansion by $W_i^\dagger$ gives the original retained
    block-diagonal matrix because $W_i W_i^\dagger=I$.  On every sector,
    \begin{align}
        \tr_{\mathrm{mult}}
        \left(\frac{\mu_\gamma}{m_\gamma}\otimes X_\gamma\right)
        &=X_\gamma,
        &
        \tr_{\mathrm{mult}}
        \left(\frac{\nu_{\sigma(\gamma)}}{m_\gamma}\otimes Y_\gamma\right)
        &=Y_\gamma,
        \notag
    \end{align}
    where the second equality uses
    $\tr(\nu_{\sigma(\gamma)})=m_\gamma$.  The two coordinate relabellings in
    the two-site composite cancel.  These identities prove
    (\ref{eq:mpdo_bnt_ambient_sector_retractions}).
\end{proof}

\begin{theorem}[Action on the weighted tensor letters]
    \label{thm:mpdo_bnt_ambient_sector_tensor_letters}
    \lean{MPOTensor.BNTAlgebraTensorClause.oneSiteAmbientSectorRetraction_verticalTensor}
    \lean{MPOTensor.BNTAlgebraTensorClause.%
        oneSiteNormalizedAmbientSectorRestoration_trace_smul_verticalTensor}
    \lean{MPOTensor.BNTAlgebraTensorClause.TwoSiteMultiplicitySpectrum.%
        twoSiteAmbientSectorRetraction_verticalTensor}
    \lean{MPOTensor.BNTAlgebraTensorClause.TwoSiteMultiplicitySpectrum.%
        twoSiteNormalizedAmbientSectorRestoration_trace_smul_verticalTensor}
    \lean{MPOTensor.verticalSectorAmbientRetraction_weightedBlockDiagonal}
    \lean{MPOTensor.normalizedVerticalSectorAmbientRestoration_trace_smul}
    \leanok
    \uses{def:mpdo_bnt_ambient_sector_maps,
        lem:mpdo_normalized_retained_vertical_sector_embedding,
        lem:mpdo_retained_vertical_sector_partial_trace,
        thm:mpdo_bnt_multiplicity_trace_normalization}
    For every vertical letter $v$, the one-site maps satisfy
    \begin{align}
        R_1(\widetilde M_v)
        &=\bigl(m_\gamma M_\gamma^v\bigr)_\gamma,
        &
        \widetilde R_1\bigl((m_\gamma M_\gamma^v)_\gamma\bigr)
        &=\widetilde M_v.
        \label{eq:mpdo_bnt_ambient_sector_one_site_letters}
    \end{align}
    The relabelled two-site maps satisfy
    \begin{align}
        R_2\!\left(\kappa_\ast(\widetilde{M^{[2]}}_v)\right)
        &=\bigl(m_\gamma A_{\sigma(\gamma)}^v\bigr)_\gamma,
        \notag\\
        \widetilde R_2\bigl((m_\gamma A_{\sigma(\gamma)}^v)_\gamma\bigr)
        &=\kappa_\ast(\widetilde{M^{[2]}}_v).
        \label{eq:mpdo_bnt_ambient_sector_two_site_letters}
    \end{align}
    The first and second one-site identities combine the one-site weighted
    sector form of
    \cite[Appendix~C.4, lines~2048--2051]{Cirac2016MPDO_arXiv} with,
    respectively,
    \cite[Appendix~C.4, lines~2067--2068]{Cirac2016MPDO_arXiv} and
    \cite[Appendix~C.4, lines~2081--2083]{Cirac2016MPDO_arXiv}.
    The first two-site identity combines the matched weighted sector form of
    \cite[Appendix~C.4, lines~2048--2064]{Cirac2016MPDO_arXiv} with the map
    $R_2$ of \cite[Appendix~C.4, lines~2078--2079]{Cirac2016MPDO_arXiv}.
    The second combines the same sector form with the normalized map
    $\widetilde R_2$ of
    \cite[Appendix~C.4, lines~2070--2071]{Cirac2016MPDO_arXiv}.
\end{theorem}

\begin{proof}\leanok
    \uses{lem:mpdo_normalized_retained_vertical_sector_embedding,
        lem:mpdo_retained_vertical_sector_partial_trace,
        thm:mpdo_bnt_multiplicity_trace_normalization}
    The vertical reconstruction identities give
    \begin{align}
        \widetilde M_v
        &=W_1^\dagger\left(\bigoplus_\gamma
          \mu_\gamma\otimes M_\gamma^v\right)W_1,
        \notag\\
        \widetilde{M^{[2]}}_v
        &=W_2^\dagger\left(\bigoplus_\gamma
          \nu_{\sigma(\gamma)}\otimes A_{\sigma(\gamma)}^v\right)W_2.
        \notag
    \end{align}
    Compression removes the two outer coisometries, and partial trace over the
    multiplicity coordinate multiplies the simple-sector letter by
    $\tr(\mu_\gamma)=m_\gamma$ or
    $\tr(\nu_{\sigma(\gamma)})=m_\gamma$.  This proves the first identity in
    each of (\ref{eq:mpdo_bnt_ambient_sector_one_site_letters}) and
    (\ref{eq:mpdo_bnt_ambient_sector_two_site_letters}).  Preparing the
    normalized multiplicity matrix cancels the factor $m_\gamma$ and the
    reconstruction identities give the two reverse formulas.
\end{proof}
